\documentclass[11pt]{article}

\usepackage[margin=1in]{geometry}
\usepackage{amsmath,amssymb,amsthm}
\usepackage{booktabs}
\usepackage{graphicx}
\usepackage{xcolor}
\usepackage{hyperref}

\newtheorem{proposition}{Proposition}
\newtheorem{remark}{Remark}

\newcommand{\Aset}{\mathcal A}
\newcommand{\Bset}{\mathcal B}
\newcommand{\Lset}{\mathcal L}
\newcommand{\Sset}{\mathcal S}
\newcommand{\Uset}{\mathcal U}
\newcommand{\Cset}{\mathcal C}
\newcommand{\RMP}{\mathrm{RMP}}
\newcommand{\LB}{\mathrm{LB}}
\newcommand{\UB}{\mathrm{UB}}
\newcommand{\eps}{\varepsilon}
\newcommand{\phidis}{\Phi^{\mathrm{dis}}}
\newcommand{\psinor}{\Psi^{\mathrm{nor}}}

\title{Certificate-Preserving Benders Acceleration for the Two-Regime DRO EVCS Model}
\author{Internal Mathematical Note}
\date{May 5, 2026}

\begin{document}
\maketitle

\section{Purpose}

This note records the mathematical algorithmic modifications added on top of
the baseline two-regime DRO EVCS planning model.  The changes are not changes
to the planning model itself.  They are changes to the Benders implementation
used to solve the disaster-resilience DRO substructure more reliably for larger
outage budgets $K$.

The important design principle is:
\begin{quote}
\emph{Intermediate acceleration may use stabilized or active-set trial points
to generate valid cuts, but convergence is claimed only after a canonical
restricted master solution passes a full-support MILP separation audit.}
\end{quote}

This separates cut-generation heuristics from the final certificate.

\section{Baseline Benders Form}

Let
\[
    x := (z,n^{\mathrm{sl}},n^{\mathrm{fa}})
\]
denote the first-stage EVCS siting and sizing vector.  A disaster support
element is denoted by
\[
    s=(b,\delta), \qquad
    b\in\Bset,\quad
    \delta\in\Uset_K
    :=\left\{\delta\in\{0,1\}^{|\Lset|}:
    \sum_{\ell\in\Lset}\delta_\ell\le K\right\}.
\]
Thus the full disaster support is
\[
    \Sset_K = \Bset\times \Uset_K .
\]

For a fixed support element $s$, let $Q_s(x)$ be the second-stage disaster
recourse value.  The DRO master contains dual-like ambiguity variables
\[
    \rho := (\alpha,\lambda)
\]
and an affine support-side expression $H_s(\rho)$.  Abstractly, the exact
master contains infinitely or combinatorially many inequalities of the form
\begin{equation}
    H_s(\rho) \ge Q_s(x), \qquad \forall s\in\Sset_K .
    \label{eq:exact-master}
\end{equation}
The restricted master keeps only a finite cut set $\Cset_t$:
\begin{equation}
    (x_t,\rho_t)\in
    \arg\min_{x,\rho}
    \left\{
    F^{\mathrm{cons}}(x)
    +(1-\pi_f)\Psi^{\mathrm{nor}}(x)
    +\pi_f D(\rho)
    :
    H_s(\rho)\ge \widehat Q_s(x),\ s\in\Cset_t
    \right\}.
    \label{eq:rmp}
\end{equation}

The separation oracle solves
\begin{equation}
    V(x_t,\rho_t)
    :=
    \max_{s\in\Sset_K}
    \left\{Q_s(x_t)-H_s(\rho_t)\right\}.
    \label{eq:full-separation}
\end{equation}
If $V(x_t,\rho_t)\le\eps$, the current canonical master solution is
$\eps$-certified.  Otherwise, the oracle returns a violated support element and
a recourse dual solution from which a valid Benders cut is generated.

\section{Cut Validity}

For a fixed support element $s$, the disaster recourse LP has a dual
representation of the form
\begin{equation}
    Q_s(x)
    =
    \max_{\pi\in\Pi_s}
    \left\{
    a_s^\top \pi + x^\top B_s^\top\pi
    \right\},
    \label{eq:recourse-dual}
\end{equation}
where $\Pi_s$ is the feasible dual polyhedron for support $s$.  Therefore, for
any dual feasible $\widehat\pi\in\Pi_s$,
\begin{equation}
    Q_s(x)\ge
    a_s^\top \widehat\pi + x^\top B_s^\top\widehat\pi
    \quad \forall x .
\end{equation}
The corresponding master cut is
\begin{equation}
    H_s(\rho)
    \ge
    a_s^\top \widehat\pi + x^\top B_s^\top\widehat\pi .
    \label{eq:benders-cut}
\end{equation}

\begin{proposition}[Intermediate Cuts Are Safe]
Any cut of the form \eqref{eq:benders-cut} generated from a valid support
element $s\in\Sset_K$ and a dual feasible vector $\widehat\pi\in\Pi_s$ is valid
for the exact master \eqref{eq:exact-master}, regardless of how $s$ was found.
\end{proposition}

\begin{proof}
For all $x$, dual feasibility gives
$a_s^\top \widehat\pi+x^\top B_s^\top\widehat\pi\le Q_s(x)$.  Since the exact
master requires $H_s(\rho)\ge Q_s(x)$, it also implies
$H_s(\rho)\ge a_s^\top \widehat\pi+x^\top B_s^\top\widehat\pi$.  Hence adding
\eqref{eq:benders-cut} cannot remove any feasible solution of the exact master.
\end{proof}

This proposition is the basis for the acceleration strategy: active sets,
stabilized trial points, or target-face tie-breaking may be used to find useful
valid cuts, but they cannot by themselves certify convergence.

\section{Algorithmic Modifications}

\subsection{Auxiliary Stabilized Trial Master}

The baseline restricted master may move to first-stage plans that are poorly
protected by the current cut pool, causing large oscillations.  To reduce this,
we added an auxiliary local-branch trial master.  Let $\bar x$ be a center plan.
The canonical master \eqref{eq:rmp} is still solved and remains the only source
of the lower bound.  Separately, the auxiliary trial master solves the same
restricted master with local restrictions such as
\begin{align}
    \sum_i |z_i-\bar z_i| &\le r_z, \label{eq:z-trust}\\
    \sum_i |n^{\mathrm{sl}}_i-\bar n^{\mathrm{sl}}_i| &\le r_{\mathrm{sl}},
    \label{eq:sl-trust}\\
    \sum_i |n^{\mathrm{fa}}_i-\bar n^{\mathrm{fa}}_i| &\le r_{\mathrm{fa}}.
    \label{eq:fa-trust}
\end{align}
The resulting trial point $x_t^{\mathrm{tr}}$ is used only for separation and
cut generation.  Its auxiliary objective value is not reported as a lower
bound.

\subsection{Active-Outage-Set Valid Cut Generation}

The full outage set $\Uset_K$ grows rapidly with $K$.  We maintain a small pool
$\mathcal P_t\subseteq \Uset_K$ containing recently active outage patterns and
their local neighborhoods, e.g.,
\[
    \mathcal N(\delta)
    =
    \{\delta': d_H(\delta,\delta')\le r,\;
    \mathbf 1^\top \delta'\le K\}.
\]
For each candidate $(b,\delta')\in\Bset\times\mathcal P_t$, the algorithm may
solve the fixed-support recourse dual and add \eqref{eq:benders-cut} if the
resulting violation is positive.  These cuts are valid because
$\Bset\times\mathcal P_t\subseteq\Sset_K$.

The active set is therefore a cut generator, not a certifier.  Certification
still requires \eqref{eq:full-separation} over the full $\Sset_K$.

\subsection{Target-Face Bundle Cuts}

When separation returns a support element $s$ and the dual recourse problem is
degenerate, an arbitrary optimal dual solution may give a weak cut.  We added a
fixed-support target-face tie-break.  At a source plan $x^0$, first compute
\[
    q^0_s := Q_s(x^0).
\]
Then restrict the dual solution to the optimal face at $x^0$:
\[
    \mathcal F_s(x^0)
    :=
    \left\{
    \pi\in\Pi_s:
    a_s^\top\pi + (x^0)^\top B_s^\top\pi
    \ge q^0_s-\tau
    \right\},
\]
where $\tau$ is a small numerical tolerance.  A target-face cut is generated by
solving
\begin{equation}
    \widehat\pi^{\mathrm{tf}}
    \in
    \arg\max_{\pi\in\mathcal F_s(x^0)}
    \left\{
    a_s^\top\pi + (\widetilde x)^\top B_s^\top\pi
    \right\},
    \label{eq:target-face}
\end{equation}
where $\widetilde x$ is a target point or core-like reference plan.  Since
$\mathcal F_s(x^0)\subseteq\Pi_s$, the generated cut is still of the valid form
\eqref{eq:benders-cut}.  The purpose is not to change the model, but to select
a dual solution that is more informative away from the source point.

\subsection{Canonical Certificate Audit}

The final implemented workflow has two phases:
\begin{enumerate}
    \item \textbf{Cut-generation phase.}  Use stabilized trial points,
    active-set cuts, and target-face bundle cuts to build a stronger cut pool.
    If a stabilized trial point reaches $V(x,\rho)\le\eps$, the row is still
    labeled diagnostic unless it is reproduced by a canonical audit.
    \item \textbf{Canonical audit phase.}  Resolve the unrestricted canonical
    restricted master with the generated cut pool, disable stabilization, and
    run full-support MILP separation.  Only if the optimal separation bound is
    below $\eps$ is the row labeled $\eps$-certified.
\end{enumerate}

This prevents auxiliary trial objectives from being confused with valid lower
bounds and keeps the convergence claim tied to the original DRO master.

\section{What Changed in the Code}

The implementation adds explicit switches for:
\begin{itemize}
    \item live iteration tracing and certificate diagnostics;
    \item warm starts and same-support cut-pool loading;
    \item local-branch trial masters on $z$ and charger counts;
    \item active outage pools and neighborhood cut generation;
    \item target-face bundle cut strengthening;
    \item canonical audit runs that disable stabilization before certification.
\end{itemize}
The EVCS model, common evaluator, and final full-support separation oracle are
not relaxed.

\section{Observed Effect on the K=10 Stress Row}

The accepted default instance uses $|\Aset|=10$, $|\Bset|=10$, and $K=10$ for
this algorithmic stress comparison.  The objective multipliers are
\[
    (m_{\mathrm{cons}},m_{\mathrm{normal}},m_{\mathrm{dis}})
    =(0.0152,1.0,1.4).
\]

\begin{table}[h]
\centering
\caption{Effect of certificate-preserving acceleration for $K=10$.}
\small
\begin{tabular}{lrrrrl}
\toprule
Method & Iter. & Cuts & Runtime (s) & Final/best violation & Status \\
\midrule
Baseline Benders & 20 & 20 & 212.40 & 107030.75 & Diagnostic \\
Local-branch stabilized & 20 & -- & 220.35 & 28095.28 & Diagnostic \\
Active-set valid cuts & 20 & 76 & 225.49 & 16861.29 & Diagnostic \\
Stabilized + active-set + target-face & 145 & 1016 & 10125.63 & 46.997 & Generator run \\
Canonical full-support audit & 1 & 1015 & 130.25 & 46.997 & $\eps$-certified \\
\bottomrule
\end{tabular}
\label{tab:k10-acceleration}
\end{table}

\begin{figure}[h]
\centering
\includegraphics[width=\linewidth]{../../results/paper_final/figures/k10_algorithm_acceleration_before_after.png}
\caption{Before--after comparison for the $K=10$ stress row.  The baseline
Benders implementation remains far above the certification tolerance.  The
combined stabilized trial, active-set valid cuts, and target-face bundle cuts
drive the full-support violation below $\eps=100$; the subsequent canonical
audit reproduces the same violation without relying on the auxiliary stabilized
objective.}
\label{fig:k10-acceleration}
\end{figure}

\section{Claim Boundary}

The certified result is currently $A=10$, $B=10$, $K=10$ for the accepted
default regime.  A larger diagnostic probe with $K=16$ is executable and
reduced the violation to $1895.79$ after 120 iterations, but it is not
certified.  Therefore, $K=16$ should be used only to diagnose the remaining
cut-strength and global convergence bottleneck, not as a paper convergence
claim.

\end{document}
